\section{How It Works}
\label{sec:method}

\subsection{Swapping boxes}
\label{sec:boxes}

A \emph{box} is one or more axes that sit next to each other. We treat a box as
one unit. Their sizes multiply together, so a box acts like a single axis of
that size.

A \emph{swap} exchanges two boxes that sit next to each other. Any reordering
of axes can be done as a list of swaps.

Here is why that helps. Cut the axis list into four parts: a front part, a left
box, a right box, and a back part. Call their sizes $D_{pre}$,
$\mathit{rows}$, $\mathit{cols}$ and $D_{post}$. Swapping the two middle boxes
is the same as transposing a $\mathit{rows} \times \mathit{cols}$ matrix. We do
it once for each of the $D_{pre}$ front positions. Each cell of that matrix is
a run of $D_{post}$ values that sit next to each other.

% Figure (fig:boxes). Body: _generated/boxes-fig.tex, written by
% pipeline/figures/plot_box_split.py.
%
% The marks are generated TikZ, drawn by this document's own build, so the
% labels are set in the paper's fonts at the sizes the generator asked for.
% Nothing here scales the picture: a \resizebox around a diagram is what makes
% one paper's figures come out with labels at four different sizes.
\begin{figure*}[!t]
\centering
% GENERATED by pipeline/figures/plot_box_split.py -- do not edit by hand.
% Hand-laid diagram; the worked shape matches plot_search_graph.py.
\begin{tikzpicture}[x=1pt,y=1pt]
\definecolor{dink}{HTML}{1C1B18}
\definecolor{dpale}{HTML}{E8EEF5}
\definecolor{dlight}{HTML}{C5D6E8}
\definecolor{dmid}{HTML}{8FB0D0}
\definecolor{ddeep}{HTML}{4A6D91}
\definecolor{daccent}{HTML}{D96B2B}
\definecolor{daccentpale}{HTML}{F7DFCE}
\useasboundingbox (0,0) rectangle (241.0,118.0);
\node[anchor=east,text=dink,align=center,font=\fontsize{5.5}{6.5}\selectfont] at (24.00,99.50) {before};
\draw[fill=dlight,draw=dink,line width=0.6pt,rounded corners=1.2pt] (30.00,92.00) rectangle (64.00,107.00);
\node[text=dink,align=center,font=\fontsize{6}{7}\selectfont] at (47.00,99.50) {$a_0$\,\fontsize{6}{7}\selectfont$=\!2$};
\draw[fill=dmid,draw=dink,line width=0.6pt,rounded corners=1.2pt] (64.00,92.00) rectangle (98.00,107.00);
\node[text=dink,align=center,font=\fontsize{6}{7}\selectfont] at (81.00,99.50) {$a_1$\,\fontsize{6}{7}\selectfont$=\!2$};
\draw[fill=dpale,draw=dink,line width=0.6pt,rounded corners=1.2pt] (98.00,92.00) rectangle (132.00,107.00);
\node[text=dink,align=center,font=\fontsize{6}{7}\selectfont] at (115.00,99.50) {$a_2$\,\fontsize{6}{7}\selectfont$=\!64$};
\draw[decorate,decoration={brace,amplitude=2.5pt,mirror},line width=0.4pt,draw=dink] (30.00,90.50) -- (64.00,90.50);
\node[anchor=north,text=dink,align=center,font=\fontsize{5.5}{6.5}\selectfont] at (47.00,87.50) {$rows$};
\draw[decorate,decoration={brace,amplitude=2.5pt,mirror},line width=0.4pt,draw=dink] (64.00,90.50) -- (98.00,90.50);
\node[anchor=north,text=dink,align=center,font=\fontsize{5.5}{6.5}\selectfont] at (81.00,87.50) {$cols$};
\draw[decorate,decoration={brace,amplitude=2.5pt,mirror},line width=0.4pt,draw=dink] (98.00,90.50) -- (132.00,90.50);
\node[anchor=north,text=dink,align=center,font=\fontsize{5.5}{6.5}\selectfont] at (115.00,87.50) {$D_{post}\!=\!64$};
\draw[-{Stealth[length=3pt,width=2.4pt]},line width=0.7pt,draw=dink] (47.00,78.00) to[bend left=28] (81.00,49.00);
\draw[-{Stealth[length=3pt,width=2.4pt]},line width=0.7pt,draw=dink] (81.00,78.00) to[bend right=28] (47.00,49.00);
\node[anchor=center,text=dink,align=center,font=\fontsize{5.5}{6.5}\selectfont] at (160.00,63.00) {the two middle boxes are\\one $rows\!\times\!cols$ matrix,\\transposed in place};
\node[anchor=east,text=dink,align=center,font=\fontsize{5.5}{6.5}\selectfont] at (24.00,37.50) {after};
\draw[fill=dlight,draw=dink,line width=0.6pt,rounded corners=1.2pt] (30.00,30.00) rectangle (64.00,45.00);
\node[text=dink,align=center,font=\fontsize{6}{7}\selectfont] at (47.00,37.50) {$a_1$\,\fontsize{6}{7}\selectfont$=\!2$};
\draw[fill=dmid,draw=dink,line width=0.6pt,rounded corners=1.2pt] (64.00,30.00) rectangle (98.00,45.00);
\node[text=dink,align=center,font=\fontsize{6}{7}\selectfont] at (81.00,37.50) {$a_0$\,\fontsize{6}{7}\selectfont$=\!2$};
\draw[fill=dpale,draw=dink,line width=0.6pt,rounded corners=1.2pt] (98.00,30.00) rectangle (132.00,45.00);
\node[text=dink,align=center,font=\fontsize{6}{7}\selectfont] at (115.00,37.50) {$a_2$\,\fontsize{6}{7}\selectfont$=\!64$};
\draw[fill=daccentpale,draw=dink,line width=0.4pt,rounded corners=0.4pt] (162.00,34.00) rectangle (170.00,42.00);
\draw[fill=white,draw=dink,line width=0.4pt,rounded corners=0.4pt] (170.00,34.00) rectangle (178.00,42.00);
\draw[fill=white,draw=dink,line width=0.4pt,rounded corners=0.4pt] (162.00,26.00) rectangle (170.00,34.00);
\draw[fill=daccentpale,draw=dink,line width=0.4pt,rounded corners=0.4pt] (170.00,26.00) rectangle (178.00,34.00);
\node[anchor=north,text=dink,align=center,font=\fontsize{5.5}{6.5}\selectfont] at (170.00,23.00) {each cell is a run of\\$D_{post}\!=\!64$ values};
\end{tikzpicture}
%
\caption{\capheadline{How one swap cuts a shape into four boxes.}
(a) The axes we have, and the order we want; shaded axes are the ones the
permutation moves. (b) The four boxes a single swap cuts them into. A box is
one or more neighbouring axes, so its extent is their product --- here the
right box covers two axes. (c) After the swap: the two middle boxes have traded
places and the outer two have not moved. Read together, the middle pair is one
$rows \times cols$ matrix, and each of its cells is a run of $D_{post}$ values
that sit next to each other in memory.}
\label{fig:boxes}
\end{figure*}


Figure~\ref{fig:boxes} shows this on a small shape. This idea is not ours. It
is the $\rho(i,j,k)$ operation from
ITTPD~\cite{cheng2025ittpd}. Section~\ref{sec:related} says more. What is ours
is how we run a swap, and how we choose which swaps to make.

\subsection{Doing a swap with no spare room}

You cannot move a value straight to where it belongs. Something else is already
there. So you pick a value up. You put it where it goes. Now you are holding
whatever was in that spot. You keep going. In the end you come back to where
you started.

That loop is called a \emph{cycle}. Walking it is called cycle
following~\cite{gustavson2012cycleleader,dummit2003abstract}. The only thing
you store is the one value in your hand. On a GPU that is one register per
thread.

In two dimensions this does not spread well over a GPU. You can only run as
many cycles at once as there are cycles.
Catanzaro~et~al.~\cite{catanzaro2014decomposition} say this plainly. A shape
with nine cycles gets nine-way parallelism, even on a machine with hundreds of
processors.

With more dimensions that limit goes away. We repeat the swap $D_{pre}$ times.
Each value we move is $D_{post}$ values wide. So the work splits into
$D_{pre} \times \mathit{cycles} \times D_{post}$ pieces that can run at the
same time. Those two extra numbers come from the axes around the swap. A
two-dimensional method does not have them.

\subsection{Skipping cells that are already correct}
\label{sec:fixedpoints}

Some cells are already in the right place. A transpose never moves the diagonal
of a square matrix.

\textbf{We skip those cells. A copy cannot.} \texttt{.contiguous()} is building
a new tensor. It must write every position, even the ones that were already
correct.

There is a formula for how many cells stay put. The transpose sends position
$p$ to $(p \bmod \mathit{cols})\,\mathit{rows} + \lfloor p/\mathit{cols}\rfloor$.
In the middle of the array this is the same as multiplying by $\mathit{cols}$
modulo $M = \mathit{rows}\cdot\mathit{cols}-1$. A cell stays put when
$p(\mathit{cols}-1) \equiv 0 \pmod M$. That has
$\gcd(\mathit{rows}-1,\mathit{cols}-1)$ answers. The last position also stays
put. So the share of the tensor a swap touches is

\begin{equation}
\mathit{moved} = 1 - \frac{\gcd(\mathit{rows}-1,\,\mathit{cols}-1) + 1}
                         {\mathit{rows}\cdot\mathit{cols}} .
\label{eq:moved}
\end{equation}

That is our whole cost model. A copy reads every value and writes every value.
It always moves $2N$ values' worth of traffic. In the cases where our method is
worth using, both methods move data at the same speed. So the speed cancels
out. Only the amount is left. That makes Equation~(\ref{eq:moved}) our
predicted time as a share of the copy's time. No number in it was measured on a
GPU. If it is below one, we do less work than the copy.

% Figure (fig:weight). Body: _generated/weight-fig.tex, written by
% pipeline/figures/plot_edge_weight.py. The fixed-cell counts drawn are
% asserted against gcd(rows-1, cols-1)+1 inside the generator, so a wrong
% picture fails `make defs` rather than printing.
\begin{figure}[!t]
\centering
\begin{tikzpicture}
\node[inner sep=2pt] (figweight) {% GENERATED by pipeline/figures/plot_edge_weight.py -- do not edit by hand.
% Fixed-cell counts recomputed from gcd(rows-1, cols-1)+1; matches src/_plan.py.
\begin{tikzpicture}[x=1pt,y=1pt]
\definecolor{dink}{HTML}{16202A}
\definecolor{dframe}{HTML}{3E4C59}
\definecolor{dground}{HTML}{F4F8FB}
\definecolor{dpre}{HTML}{DCE9F2}
\definecolor{dleft}{HTML}{A9CFE3}
\definecolor{dright}{HTML}{6FA8C9}
\definecolor{dpost}{HTML}{CFE0EC}
\definecolor{dband}{HTML}{2E5A75}
\definecolor{daccent}{HTML}{C2571A}
\definecolor{daccentpale}{HTML}{F6E2D2}
\useasboundingbox (0,0) rectangle (241.0,120.0);
\draw[draw=dframe,line width=0.7pt,rounded corners=2.4pt] (0,0) rectangle (241.0,120.0);
\draw[fill=dground,draw=dframe,line width=0.5pt,rounded corners=1.6pt] (6.00,6.00) rectangle (118.50,114.00);
\node[anchor=north west,text=dink,align=center,font=\fontsize{6}{7}\selectfont\bfseries] at (10.00,110.00) {(a)};
\draw[fill=daccentpale,draw=dink,line width=0.35pt] (42.25,76.00) rectangle (62.25,96.00);
\node[anchor=center,text=daccent,align=center,font=\fontsize{4.6}{5.4}\selectfont] at (52.25,86.00) {\textbf{0}};
\draw[fill=white,draw=dink,line width=0.35pt] (62.25,76.00) rectangle (82.25,96.00);
\node[anchor=center,text=dink,align=center,font=\fontsize{4.6}{5.4}\selectfont] at (72.25,86.00) {1};
\draw[fill=white,draw=dink,line width=0.35pt] (42.25,56.00) rectangle (62.25,76.00);
\node[anchor=center,text=dink,align=center,font=\fontsize{4.6}{5.4}\selectfont] at (52.25,66.00) {2};
\draw[fill=daccentpale,draw=dink,line width=0.35pt] (62.25,56.00) rectangle (82.25,76.00);
\node[anchor=center,text=daccent,align=center,font=\fontsize{4.6}{5.4}\selectfont] at (72.25,66.00) {\textbf{3}};
\node[anchor=center,text=dink,align=center,font=\fontsize{6}{7}\selectfont] at (62.25,105.00) {a $2\!\times\!2$ block};
\draw[fill=dpre,draw=dink,line width=0.5pt,rounded corners=1.2pt] (11.00,11.00) rectangle (113.50,35.00);
\node[anchor=center,text=dink,align=center,font=\fontsize{4.6}{5.4}\selectfont] at (62.25,28.50) {$\gcd(2\!-\!1,2\!-\!1)+1 = 2$ stay put};
\node[anchor=center,text=dink,align=center,font=\fontsize{5.2}{6.2}\selectfont] at (62.25,18.00) {$moved = 1-\tfrac{2}{4} = \mathbf{0.50}$};
\draw[fill=dground,draw=dframe,line width=0.5pt,rounded corners=1.6pt] (122.50,6.00) rectangle (235.00,114.00);
\node[anchor=north west,text=dink,align=center,font=\fontsize{6}{7}\selectfont\bfseries] at (126.50,110.00) {(b)};
\draw[fill=daccentpale,draw=dink,line width=0.35pt] (150.75,82.00) rectangle (164.75,96.00);
\node[anchor=center,text=daccent,align=center,font=\fontsize{4.6}{5.4}\selectfont] at (157.75,89.00) {\textbf{0}};
\draw[fill=white,draw=dink,line width=0.35pt] (164.75,82.00) rectangle (178.75,96.00);
\node[anchor=center,text=dink,align=center,font=\fontsize{4.6}{5.4}\selectfont] at (171.75,89.00) {1};
\draw[fill=white,draw=dink,line width=0.35pt] (178.75,82.00) rectangle (192.75,96.00);
\node[anchor=center,text=dink,align=center,font=\fontsize{4.6}{5.4}\selectfont] at (185.75,89.00) {2};
\draw[fill=white,draw=dink,line width=0.35pt] (192.75,82.00) rectangle (206.75,96.00);
\node[anchor=center,text=dink,align=center,font=\fontsize{4.6}{5.4}\selectfont] at (199.75,89.00) {3};
\draw[fill=white,draw=dink,line width=0.35pt] (150.75,68.00) rectangle (164.75,82.00);
\node[anchor=center,text=dink,align=center,font=\fontsize{4.6}{5.4}\selectfont] at (157.75,75.00) {4};
\draw[fill=daccentpale,draw=dink,line width=0.35pt] (164.75,68.00) rectangle (178.75,82.00);
\node[anchor=center,text=daccent,align=center,font=\fontsize{4.6}{5.4}\selectfont] at (171.75,75.00) {\textbf{5}};
\draw[fill=white,draw=dink,line width=0.35pt] (178.75,68.00) rectangle (192.75,82.00);
\node[anchor=center,text=dink,align=center,font=\fontsize{4.6}{5.4}\selectfont] at (185.75,75.00) {6};
\draw[fill=white,draw=dink,line width=0.35pt] (192.75,68.00) rectangle (206.75,82.00);
\node[anchor=center,text=dink,align=center,font=\fontsize{4.6}{5.4}\selectfont] at (199.75,75.00) {7};
\draw[fill=white,draw=dink,line width=0.35pt] (150.75,54.00) rectangle (164.75,68.00);
\node[anchor=center,text=dink,align=center,font=\fontsize{4.6}{5.4}\selectfont] at (157.75,61.00) {8};
\draw[fill=white,draw=dink,line width=0.35pt] (164.75,54.00) rectangle (178.75,68.00);
\node[anchor=center,text=dink,align=center,font=\fontsize{4.6}{5.4}\selectfont] at (171.75,61.00) {9};
\draw[fill=daccentpale,draw=dink,line width=0.35pt] (178.75,54.00) rectangle (192.75,68.00);
\node[anchor=center,text=daccent,align=center,font=\fontsize{4.6}{5.4}\selectfont] at (185.75,61.00) {\textbf{10}};
\draw[fill=white,draw=dink,line width=0.35pt] (192.75,54.00) rectangle (206.75,68.00);
\node[anchor=center,text=dink,align=center,font=\fontsize{4.6}{5.4}\selectfont] at (199.75,61.00) {11};
\draw[fill=white,draw=dink,line width=0.35pt] (150.75,40.00) rectangle (164.75,54.00);
\node[anchor=center,text=dink,align=center,font=\fontsize{4.6}{5.4}\selectfont] at (157.75,47.00) {12};
\draw[fill=white,draw=dink,line width=0.35pt] (164.75,40.00) rectangle (178.75,54.00);
\node[anchor=center,text=dink,align=center,font=\fontsize{4.6}{5.4}\selectfont] at (171.75,47.00) {13};
\draw[fill=white,draw=dink,line width=0.35pt] (178.75,40.00) rectangle (192.75,54.00);
\node[anchor=center,text=dink,align=center,font=\fontsize{4.6}{5.4}\selectfont] at (185.75,47.00) {14};
\draw[fill=daccentpale,draw=dink,line width=0.35pt] (192.75,40.00) rectangle (206.75,54.00);
\node[anchor=center,text=daccent,align=center,font=\fontsize{4.6}{5.4}\selectfont] at (199.75,47.00) {\textbf{15}};
\node[anchor=center,text=dink,align=center,font=\fontsize{6}{7}\selectfont] at (178.75,105.00) {a $4\!\times\!4$ block};
\draw[fill=dpre,draw=dink,line width=0.5pt,rounded corners=1.2pt] (127.50,11.00) rectangle (230.00,35.00);
\node[anchor=center,text=dink,align=center,font=\fontsize{4.6}{5.4}\selectfont] at (178.75,28.50) {$\gcd(4\!-\!1,4\!-\!1)+1 = 4$ stay put};
\node[anchor=center,text=dink,align=center,font=\fontsize{5.2}{6.2}\selectfont] at (178.75,18.00) {$moved = 1-\tfrac{4}{16} = \mathbf{0.75}$};
\end{tikzpicture}
};
\draw[black,line width=0.5pt,rounded corners=2pt]
  (figweight.south west) rectangle (figweight.north east);
\end{tikzpicture}
\caption{\capheadline{Where an edge's weight comes from.}
Shaded cells are already in the right place. A transpose does not move them, so
we do not either, and a copy has to write them anyway. The smaller the block,
the larger the share that stays put. A
$\modelMinMovedRows \times \modelMinMovedCols$ block is the best any single
swap can do.}
\label{fig:weight}
\end{figure}


Figure~\ref{fig:weight} draws this for two blocks. Now the main point. Any
method that builds output positions must go over the
tensor at least once. It has to write every position. Ours does not. Swapping a
$\modelMinMovedRows \times \modelMinMovedCols$ block moves $\modelMinMoved$ of
the tensor. No single swap can do better. The reason is that
$\gcd(\mathit{rows}-1,\mathit{cols}-1)$ is at most
$\min(\mathit{rows},\mathit{cols})-1$, so the skipped share is at most
$1/\max(\mathit{rows},\mathit{cols})$.

Two things follow. Only a method that skips can do part of a pass. And two
swaps can never add up to less than one pass. So \emph{beating
\texttt{.contiguous()} and using a single swap are the same thing}.

\subsection{Choosing which swaps to make}
\label{sec:planning}

There are many lists of swaps that reach the same order. They do not cost the
same.

We build a graph. Each way of ordering the axes is a point. Each legal swap is
a line between two points. The weight on that line is how many bytes the swap
moves. We know that number exactly from Equation~(\ref{eq:moved}). Then we run
Dijkstra's algorithm to find the cheapest route.

% Figure (fig:graph). Body: _generated/graph-fig.tex, written by
% pipeline/figures/plot_search_graph.py, whose weights come from the shipped
% planner rather than from a table typed here.
%
% figure*, not figure: six nodes laid out as the animation lays them out need
% the full two-column measure. At 241pt the node boxes are under 40pt wide and
% the edge labels collide.
\begin{figure*}[!t]
\centering
\begin{tikzpicture}
\node[inner sep=2pt] (figgraph) {% GENERATED by pipeline/figures/plot_search_graph.py -- do not edit by hand.
% Weights from src/_plan.py; layout mirrors docs/dijkstra-animation.html.
\begin{tikzpicture}[x=1pt,y=1pt]
\definecolor{dink}{HTML}{16202A}
\definecolor{dframe}{HTML}{3E4C59}
\definecolor{dground}{HTML}{F4F8FB}
\definecolor{dpre}{HTML}{DCE9F2}
\definecolor{dleft}{HTML}{A9CFE3}
\definecolor{dright}{HTML}{6FA8C9}
\definecolor{dpost}{HTML}{CFE0EC}
\definecolor{dband}{HTML}{2E5A75}
\definecolor{daccent}{HTML}{C2571A}
\definecolor{daccentpale}{HTML}{F6E2D2}
\useasboundingbox (0,0) rectangle (500.0,150.0);
\draw[draw=dframe,line width=0.7pt,rounded corners=2.4pt] (0,0) rectangle (500.0,150.0);
\draw[-{Stealth[length=3pt,width=2.4pt]},line width=0.35pt,draw=dright] (61.00,78.67) to[bend left=9] (313.83,120.83);
\draw[-{Stealth[length=3pt,width=2.4pt]},line width=0.35pt,draw=dright] (59.24,63.50) to[bend left=9] (135.44,33.00);
\draw[-{Stealth[length=3pt,width=2.4pt]},line width=0.35pt,draw=dright] (61.00,69.33) to[bend left=9] (313.83,27.17);
\draw[-{Stealth[length=3pt,width=2.4pt]},line width=0.35pt,draw=dright] (135.44,115.00) to[bend left=9] (59.24,84.50);
\draw[-{Stealth[length=3pt,width=2.4pt]},line width=0.35pt,draw=dright] (189.68,125.50) to[bend left=9] (313.83,125.50);
\draw[-{Stealth[length=3pt,width=2.4pt]},line width=0.35pt,draw=dright] (180.04,115.00) to[bend left=9] (323.47,33.00);
\draw[-{Stealth[length=3pt,width=2.4pt]},line width=0.35pt,draw=dright] (313.83,120.83) to[bend left=9] (61.00,78.67);
\draw[-{Stealth[length=3pt,width=2.4pt]},line width=0.35pt,draw=dright] (313.83,125.50) to[bend left=9] (189.68,125.50);
\draw[-{Stealth[length=3pt,width=2.4pt]},line width=0.35pt,draw=dright] (323.47,115.00) to[bend left=9] (180.04,33.00);
\draw[-{Stealth[length=3pt,width=2.4pt]},line width=0.35pt,draw=dright] (367.35,115.00) to[bend left=9] (441.48,84.50);
\draw[-{Stealth[length=3pt,width=2.4pt]},line width=0.35pt,draw=dright] (135.44,33.00) to[bend left=9] (59.24,63.50);
\draw[-{Stealth[length=3pt,width=2.4pt]},line width=0.35pt,draw=dright] (180.04,33.00) to[bend left=9] (323.47,115.00);
\draw[-{Stealth[length=3pt,width=2.4pt]},line width=0.35pt,draw=dright] (189.68,22.50) to[bend left=9] (313.83,22.50);
\draw[-{Stealth[length=3pt,width=2.4pt]},line width=0.35pt,draw=dright] (189.68,27.22) to[bend left=9] (439.00,69.28);
\draw[-{Stealth[length=3pt,width=2.4pt]},line width=0.35pt,draw=dright] (313.83,27.17) to[bend left=9] (61.00,69.33);
\draw[-{Stealth[length=3pt,width=2.4pt]},line width=0.35pt,draw=dright] (323.47,33.00) to[bend left=9] (180.04,115.00);
\draw[-{Stealth[length=3pt,width=2.4pt]},line width=0.35pt,draw=dright] (313.83,22.50) to[bend left=9] (189.68,22.50);
\draw[-{Stealth[length=3pt,width=2.4pt]},line width=0.35pt,draw=dright] (367.35,33.00) to[bend left=9] (441.48,63.50);
\draw[-{Stealth[length=3pt,width=2.4pt]},line width=0.35pt,draw=dright] (439.00,78.72) to[bend left=9] (189.68,120.78);
\draw[-{Stealth[length=3pt,width=2.4pt]},line width=0.35pt,draw=dright] (441.48,84.50) to[bend left=9] (367.35,115.00);
\draw[-{Stealth[length=3pt,width=2.4pt]},line width=0.35pt,draw=dright] (439.00,69.28) to[bend left=9] (189.68,27.22);
\draw[-{Stealth[length=3pt,width=2.4pt]},line width=0.35pt,draw=dright] (441.48,63.50) to[bend left=9] (367.35,33.00);
\draw[-{Stealth[length=3pt,width=2.4pt]},line width=1.5pt,draw=daccent] (59.24,84.50) to[bend left=9] (135.44,115.00);
\node[anchor=center,text=daccent,align=center,font=\fontsize{6}{7}\selectfont,fill=white,inner sep=1.2pt,rounded corners=0.8pt] at (85.59,102.59) {128\,B};
\draw[-{Stealth[length=3pt,width=2.4pt]},line width=1.5pt,draw=daccent] (189.68,120.78) to[bend left=9] (439.00,78.72);
\node[anchor=center,text=daccent,align=center,font=\fontsize{6}{7}\selectfont,fill=white,inner sep=1.2pt,rounded corners=0.8pt] at (285.59,111.70) {176\,B};
\draw[fill=dleft,draw=dink,line width=0.5pt] (7.00,65.50) rectangle (59.00,82.50);
\node[anchor=center,text=dink,align=center,font=\fontsize{6}{7}\selectfont] at (33.00,74.00) {(ABC)};
\node[anchor=south,text=dink,align=center,font=\fontsize{5.2}{6.2}\selectfont] at (33.00,85.50) {start};
\draw[fill=dpre,draw=dink,line width=0.5pt] (135.68,117.00) rectangle (187.68,134.00);
\node[anchor=center,text=dink,align=center,font=\fontsize{6}{7}\selectfont] at (161.68,125.50) {(BAC)};
\draw[fill=dpre,draw=dink,line width=0.5pt] (315.83,117.00) rectangle (367.83,134.00);
\node[anchor=center,text=dink,align=center,font=\fontsize{6}{7}\selectfont] at (341.83,125.50) {(BCA)};
\draw[fill=dpre,draw=dink,line width=0.5pt] (135.68,14.00) rectangle (187.68,31.00);
\node[anchor=center,text=dink,align=center,font=\fontsize{6}{7}\selectfont] at (161.68,22.50) {(CAB)};
\draw[fill=dpre,draw=dink,line width=0.5pt] (315.83,14.00) rectangle (367.83,31.00);
\node[anchor=center,text=dink,align=center,font=\fontsize{6}{7}\selectfont] at (341.83,22.50) {(ACB)};
\draw[fill=daccentpale,draw=dink,line width=0.5pt] (441.00,65.50) rectangle (493.00,82.50);
\node[anchor=center,text=dink,align=center,font=\fontsize{6}{7}\selectfont] at (467.00,74.00) {(CBA)};
\node[anchor=south,text=dink,align=center,font=\fontsize{5.2}{6.2}\selectfont] at (467.00,85.50) {target};
\node[anchor=north,text=dink,align=center,font=\fontsize{5.2}{6.2}\selectfont] at (250.00,11.00) {Shape $(2,3,4)$, target $(2,1,0)$. Every ordering is a node and every legal swap an edge. The cheapest route costs 304\,bytes.};
\end{tikzpicture}
};
\draw[black,line width=0.5pt,rounded corners=2pt]
  (figgraph.south west) rectangle (figgraph.north east);
\end{tikzpicture}
\caption{\capheadline{The graph the planner searches.}
Each node is one way of ordering the axes. Each arrow is a legal swap. The
cheapest route is drawn in orange with the bytes each of its swaps moves. The
weights are not equal, so the cheapest route is not simply the shortest one.
This is the same example, and the same layout, as the step-by-step animation in
\texttt{docs/dijkstra-animation.html}; the generator checks its cost model
against the planner's, so the two cannot drift apart.}
\label{fig:graph}
\end{figure*}


Figure~\ref{fig:graph} shows the start of that search on a small shape. We do
not remove any lines. We do not apply any cutoff to them. We only reject
a finished route. If the cheapest route needs more than $\modelMaxSteps$ swaps,
we return nothing. The caller then uses \texttt{.contiguous()}. Every swap
rewrites the whole buffer, so a long route moves several tensors' worth of data.

The search does not make us faster than a copy. Section~\ref{sec:fixedpoints}
shows that needs a single swap, and you do not need a search to find one. The
search matters for the longer routes. Nothing beats a copy there, but the
search keeps us ahead of the other in-place methods.

\subsection{Planning happens once}
\label{sec:cache}

There are three phases, and only one of them runs often.

\emph{Planning} builds the graph and runs the search. It touches no data at
all. It only needs the shape and the permutation, so it runs on the CPU.

\emph{Preparing} uploads what the kernel needs to the GPU. \emph{Applying}
runs the swaps on the real tensor. That last one is the only phase whose cost
depends on how big the tensor is.

Planning is not free. For a single swap it returns at once, because the search
stops as soon as it reaches the target. For a hard permutation at rank~8 it
took several seconds (Section~\ref{sec:limits}).

That would be a problem if we paid it every call. We do not. A plan depends
only on the shape and the permutation, and neither changes from one call to the
next. So we work it out once and keep it. The same tensor shape appears again
and again in a training loop, so the first call pays for the plan and every
call after it pays nothing.

This is why planning cost does not appear in any of our timings. We time
applying, which is what runs in the loop. We say so here so that nobody reads
the numbers in Section~\ref{sec:eval} as though the search ran inside them.

\subsection{Nothing is tuned}

An older version of this work used three numbers measured on one GPU. There was
a throughput knee, a routing cutoff, and an average bandwidth. All three are
gone. The bandwidth cancelled out. The knee became a question about the shape
instead of a measurement.

We can check that the model knows where it works. Inside the region we claim,
Equation~(\ref{eq:moved}) predicts the measured ratio to within
$\modelErrIn\%$. Outside it, the error is about $\modelErrOut\%$. That is why
the region exists.
